\chapter{121}



Wie oft sich Louis diese erste Begegnung vorgestellt hatte. Vor allem, ob er es überhaupt aushalten würde, sie in dieser ungewohnten Situation zu sehen. Ene Frau wie sie, an einen Rollstuhl gebunden wie eine Gefangene. 
Und jetzt, wie er sie über den Weg vor sich hinschob, war, als würde ihm
Zeit geschenkt, um in ihr Wiedersehen hineinzuwachsen. Denn Giorgia
redete eine geraume Zeit ohne Unterbruch. Gelegentlich schaute sie über
die Schulter schräg zu ihm hoch, als wollte sie sich vergewissern, dass
er noch da war.

Sie erzählte von ihrer Zeit im \emph{Ospedale}, ihrer masslosen Angst,
von nun an ein Krüppel zu sein, von den vielen Besuchen, und dass sie
offensichtlich erst die Letzten wach registriert hatte. Vom Schmerz
ihrer Eltern, und wie hilflos sie neben ihrem Bett gestanden hatten.
Aber auch von den Ärzten, von all den freundlichen Menschen, die sich
dort rührend um sie gekümmert hätten. Die Aufklärung des Unfalls
erwähnte sie nebenbei. Für ein paar Sätze kam sie in Fahrt.

«Dieses endlose, zermürbende Verfahren kann mir gestohlen
bleiben. Nichts kann die Sache ungeschehen machen. Aber meine
Eltern wollen Klarheit. Ich kann sie ja verstehen. Sie haben einen
Anwalt engagiert.»

Mit einer seltsam abwinkenden Geste schien sie das Thema verbannen zu
wollen.
Louis’ Verwirrung wuchs. Mit keinem Wort hatte Giorgia sich zur besagten
Nacht auf dem \emph{Roccia} geäussert. Weder zu Bastiano, zu allfälligen
Drogen noch zu ihrer Begegnung. Waren tatsächlich all ihre Erinnerungen
an diese Nacht gelöscht?
Wie sollte er den Einstieg finden, sich ihr zu erklären?
Unverhofft hellte sich Giorgias Gesicht auf.

«Eines Tages stand plötzlich Valentina an meinem Bett. Ich hatte wohl
gedöst und sie nicht eintreten hören,» Giorgia stoppte die Räder des
Rollstuhls so plötzlich, dass Louis unsanft mit dem Bauch an das Gefährt
stiess, «verzeih.»

Louis stellte sich neben sie.

«Schon gut.»

Er ging wieder vor ihr in die Hocke. Während ihr Blick auf dem Meer
liegen blieb, fuhr sie fort.

«In diesem Moment wusste ich, ich bin da, oder wie du willst, wieder da. Ein eigenartiges Gefühl,
schwer zu erklären. Die Ärzte meinten, ich sei zuvor mehrmals kurz zu
mir gekommen, aber gleich wieder abgetaucht. Das sei eine gesunde
Schutzreaktion. Der Körper brauche in solchen Situationen seine ganze
Kraft für die Heilung. Aber egal,» sie wandte sich Louis zu und ihre
Augen wurden feucht, «Valentinas Besuch war mein erstes Highlight. Ich,
ich konnte mit ihr sprechen, wie, wie immer.»

Zögernd legte Louis seine Hand auf ihren Arm. Sie liess es geschehen.

«Ja, Louis, und seit ein paar Wochen kämpfe ich mich zurück, manchmal
verdammt wütend, manchmal komplett überzeugt, dass ich wieder ganz zu
mir werde.»

Louis schaute sie lange an, bis Giorgia den Blick senkte und sich
einzelne Tränen abwischte.
Über ihnen setzten sich hellgraue Wolken in Richtung Süden in Bewegung.
Sie glitten zögernd vor die Sonne. Der Himmel verdunkelte sich. Nach und
nach wurde es kühler.
Giorgia unterbrach die Stille.

«Ohne Hilfe komme ich nicht mehr ans Wasser. Schau mal,» sie zeigte auf
eine kleine Bucht in der Ferne, «bringst du uns dahin? Das wäre schön.
Und dann bist du dran. Ich bin gespannt, was bei dir in der Zwischenzeit
gelaufen ist.»

Es war wieder eine dieser emotionalen Wendungen, die ihn jedes Mal
umhauten.

«Ja, klar, mach ich gerne.»

Die Bucht war über eine schmale Steintreppe zu erreichen. Louis schaute
sich um. Es gab keinen andern Zugang. Giorgia schien erwartet zu haben,
dass er sie hinunter trug. Sie würden sich sehr nah kommen. Ganz
selbstverständlich streckte sie die Arme nach ihm aus. 
Louis hob sie mit sicherem Griff aus dem Rollstuhl. Sie umfasste seinen
Hals und er zirkelte vorsichtig die Treppe hinunter. Dann trug er
sie über den halbrunden, kleinen Platz zu einem geeigneten
Sitzstein. Sanft setzte er sie ab, rückte ihr die Beine zurecht, holte
oben den Rollstuhl und ging schliesslich vor ihr in die Knie. Er zog
eine Flasche aus seinem Rucksack. 

«Hast du Durst? Das ist Apfelsaft aus den Schweizer Bergen.»

Giorgia nickte überrascht und als sie trank, bemerkte er, dass sie zitterte.

«Ist dir kalt?»

«Ein bisschen. Wenn ich friere, muss ich dauernd aufs Klo. Und das ist
verdammt kompliziert geworden. In der Tasche, schau unten im Rollstuhl, 
hat es eine Decke. Gibst du sie mir, bitte?»

«Natürlich.»

Während Louis die Tasche öffnete und die Decke heraushob, wusste er, die Uhr tickte. Der Countdown lief.
Endgültig. Unaufschiebbar. Es konnte sich höchstens noch um Minuten handeln, bis er es tat.
Er hob Giorgia in den Rollstuhl zurück und legte die Decke über die Vorderseite ihres Körpers. 
Auf den Seiten schob er sie dicht an ihre Beine.
Die ganze Zeit über hatten sie es vermieden, von ihnen beiden zu
sprechen. Louis fühlte sich wie als Kind, wenn er sich Nachts davor
gefürchtet hatte, unter sein Bett zu schauen. Jahrelang hatte er dort
ein Ungeheuer vermutet, das plötzlich hervorkommen und ihn auffressen
könnte. Die Uhr tickte weiter.

«Giorgia.»

Sie reckte sofort den Hals, ihre Augen vergrösserten sich und mit einem
Nicken schien sie ihn ermuntern zu wollen, weiterzusprechen.

«Ich hatte eine schreckliche Zeit. Ich hielt mich in der Schweiz
versteckt, weil ich, ich werde dir alles erzählen.»

Giorgia verschränkte die Finger ineinander, schien Luft zu holen und auf
ihrem Gesicht war ein Anflug von Furcht zu erkennen.

«Du weisst, nach unserer Trennung fiel ich zusammen wie ein Kind. Ich
habe um deine Liebe gebettelt. Es tut mir leid. Ich war ein Versager.
Seit Wochen bin ich mit einem Mentor dran, mein Leben zu sortieren. Joh
ist ein bemerkenswerter Mann. Er hat mich hier hin begleitet. Aber darum
geht es gerade nicht. Wenn du willst, erzähle ich dir ein anderes Mal
von der \emph{Sunnewaid,} also von dem Ort, von dem ich aktuell komme
und an den ich vorläufig wieder zurückkehren werde.»

Wieder nickte sie mehrmals. Sie hatte die Decke unterdessen bis zum Hals hochgezogen.

«Du hast mich immer unterstützt, wenn es um meine Musik ging, Giorgia,
immer. Dafür bin ich dir so dankbar,» sie lächelte, «ich habe dir etwas
mitgebracht.»

Louis zog einen kleinen Boomer aus dem Rucksack.

«Bevor ich dir die wirklich wichtige Sache erzähle, möchte ich dir einen
Song abspielen,» er winkte ab, «nein, keiner von mir, aber etwas ganz
besonderes. Er beschreibt sehr passend, wie es mir nach unserer Trennung
ging und eigentlich immer noch geht. Stell dir vor, die Worte seien von
mir, Giorgia. Vielleicht gefällt er dir.»

Giorgia sank tiefer in den Rollstuhl, legte den Kopf schräg und sah
zunehmend verdutzter aus.
Louis zog sein Handy aus der Hosentasche.

«Es ist von Leonard Cohen, also er lebt ja nicht mehr, aber seine Söhne
haben nach seinem Tod letzte Songs von ihm vollendet und produziert.
Bald erscheint eine neue CD. Der Song, den ich dir abspiele, heisst
\emph{«Moving on»}».

«Von Cohen?» ein Strahlen flog über ihr Gesicht, «echt? Aber wie, die
erscheint bald?»

«Ja, da gibt es einen kleinen Haken. Die CD ist noch nicht
veröffentlicht. Ich vertraue darauf, dass du niemandem davon erzählst.»

«Aber, wie kommt es, dass du mir einen unveröffentlichten Song abspielen
kannst?»

«Das darf ich dir nicht sagen, Giorgia.»

Sie wirkte überrascht, schmunzelte und fixierte das Handy in Louis’
Hand.

«Und ja, bald hätte ich es vergessen. Den Teil mit den Hormonen 
überhörst du bitte, der passt nicht.»

Er setzte sich vor Giorgia auf den Boden und stellte den Boomer zwischen
ihnen auf. Mit dem erneuten Auftauchen der Sonne klappte Giorgia
erleichtert die Decke zur Seite.
Louis startete den Song.Er schloss die Augen und stellte sich vor, der Mann
sei für eine Songlänge zurückgekehrt. Exklusiv für sie beide. 

Als sässen sie direkt neben der Kirche, begannen die Glocken zu läuten und mit 
Cohens Stimme legte sich wellenartig eine erdende Zärtlichkeit über den kleinen Ort.   
 

\qrblock{Leonard Cohen \emph{«Moving on»}}{media/QR-Codes/002_Moving_On_Leonard_Cohen_Spotify.png}

\newpage 

\songtexttitled{«Moving on»}{
Ich liebte dein Gesicht, ich liebte dein Haar,\\
Deine T-Shirts und deine Abendgarderobe.\\
Was die Welt, den Job und den Krieg angeht,\\
Ich habe sie alle weggeworfen, um dich mehr zu lieben.\\
Und jetzt bist du weg, jetzt bist du weg,\\
Als ob es jemals ein Du gegeben hätte.\\
Wer brach das Herz und machte es neu?\\
Wer macht weiter, wer verarscht hier wen?\\
Ich liebte deine Launen, ich liebe die Art,\\
Sie bedrohen jeden einzelnen Tag.\\
Deine Schönheit beherrschte mich, obwohl ich wusste\\
dass es mehr hormonell war als die Aussicht.\\
Jetzt bist du weg, jetzt bist du weg,\\
Als ob es dich jemals gegeben hätte.\\
Königin des Flieders, Königin des Blaus.\\
Wer zieht weiter, wer verarscht wen?\\
Ich liebte dein Gesicht, ich liebte dein Haar,\\
Deine T-Shirts und deine Abendgarderobe.\\
Was die Welt, den Job, den Krieg angeht,\\
Ich liess sie alle hinter mir, um dich mehr zu lieben.\\
Und jetzt bist du weg, jetzt bist du weg,\\
Als ob es jemals ein Du gegeben hätte.\\
Wer hielt mich im Sterben, zog mich durch,\\
Wer zieht weiter, wer verarscht wen?\\
Wer zieht weiter, wer verarscht hier wen?}
